\documentclass{beamer}
\usetheme{Madrid}

\title{What Happens When Python Runs?}
\author{CPSC 250}
\date{}

\begin{document}

\begin{frame}
\titlepage
\end{frame}

\begin{frame}[fragile]{The Big Picture}

When you run

\begin{verbatim}
python myprogram.py
\end{verbatim}

Python is:

\begin{itemize}
\item Not directly executed by the CPU
\item Not translated into C++ source code
\item Executed through an interpreter and virtual machine
\end{itemize}

\end{frame}

\begin{frame}[fragile]{Step 1: Read the Source File}

Example:

\begin{verbatim}
x = 5
y = 7
print(x + y)
\end{verbatim}

The Python interpreter first reads the text file.

\end{frame}

\begin{frame}[fragile]{Step 2: Parsing}

The interpreter converts the text into an Abstract Syntax Tree.

\begin{verbatim}
Assignment
 ├── Variable("x")
 └── Integer(5)
\end{verbatim}

\end{frame}

\begin{frame}[fragile]{Step 3: Compile to Bytecode}

The AST is converted into Python bytecode.

\begin{verbatim}
LOAD_CONST 5
STORE_NAME x
LOAD_CONST 7
STORE_NAME y
LOAD_NAME x
LOAD_NAME y
BINARY_OP +
CALL print
\end{verbatim}

Bytecode is lower-level than Python source,
but it is still not machine code.

\end{frame}

\begin{frame}[fragile]{Step 4: The Python Virtual Machine}

The bytecode is executed by the Python Virtual Machine.

\begin{verbatim}
while True:
    instruction = next_bytecode()
    execute(instruction)
\end{verbatim}

The virtual machine acts like a tiny computer
whose instruction set is Python bytecode.

\end{frame}

\begin{frame}[fragile]{Where Does C Come In?}

CPython is written mostly in C.

For example:

\begin{verbatim}
x + y
\end{verbatim}

may ultimately call something like:

\begin{verbatim}
PyNumber_Add(x, y)
\end{verbatim}

The CPU executes the compiled C interpreter,
not the original Python source.

\end{frame}

\begin{frame}[fragile]{Python Objects in Memory}

Everything in Python is an object.

\begin{verbatim}
x = 42
\end{verbatim}

creates an integer object, and the name
\texttt{x} refers to that object.

\begin{verbatim}
x  --->  PyLongObject(42)
\end{verbatim}

\end{frame}

\begin{frame}{The Complete Pipeline}

\begin{center}
Python Source

$\downarrow$

Parser

$\downarrow$

Abstract Syntax Tree

$\downarrow$

Bytecode Compiler

$\downarrow$

Python Bytecode

$\downarrow$

Python Virtual Machine

$\downarrow$

C Implementation of Python

$\downarrow$

Machine Code on the CPU
\end{center}

\end{frame}

\begin{frame}{Takeaway}

When you run a Python program:

\begin{enumerate}
\item Python source code is parsed
\item An AST is created
\item The AST is compiled to bytecode
\item The virtual machine executes the bytecode
\item The interpreter's compiled C code runs on the CPU
\end{enumerate}

\vspace{0.4cm}

Python is often called interpreted, but CPython
also performs a compilation step.

\end{frame}

\end{document}
